% Folder 119, batch 15 — archivists' pages 281–300.
% Pass: Opus 5.5 (claude-opus-5-5), 2026-10-10 — notes only (issue #49), written from the survey of the facsimile of 3 October 2026 (issue #45); no page of the batch is transcribed.
%
% Contents of the batch (special rules for folder 119, issues #45 and #49):
%  - pp. 281, 287: blank archivists' divider sheets. No \page.
%  - pp. 282–286: typed letter of Ronald Brown to Grothendieck, « 16.04.91 »,
%    in English, signed by hand. Third-party correspondence: under RIGHTS.md
%    it is not circulated without permission, so each page carries a note
%    (writer, date, subject) and nothing of its text. Decided 2026-10-10, as
%    for the letters of Mazur and Deligne in 144 and of Deligne in 32.
%  - pp. 288–300: Serre's Collège de France résumé 1990–91 (ends p. 301, batch
%    16). Third-party and printed: notes only.
\documentclass[11pt,a4paper]{article}
\input{../preamble/grothendieck}

\folder{119}
\batch{15}
\pages{281}{300}
\foldertitle{Esquisse d’un programme [dossier constitué en vue d'une candidature au CNRS (1984)] : copies de tapuscrits et de tapuscrits annotés (1972-1974, 1978-1979, 1984, 1990-1991, s.d.), tirés à part (1971), notes et copies de note manuscrites (1986, s.d.), lettres (1968, 1972, 1991).}
\dating{1968-1991}
\watermark{Édition de démonstration}

\begin{document}

\section*{Lettre de Ronald Brown, 16 avril 1991 (pages 282 à 286)}

\page{282}
\note{Lettre dactylographiée de Ronald Brown (Bangor) à Grothendieck, en anglais, datée « 16.04.91 », signée à la main, pp.~282--286. Elle répond à la lettre de Grothendieck du 9.4.1991 (imprimée dans l'édition de la correspondance Grothendieck--Brown, qui ne donne pas cette réponse) : ses propres recherches, la symétrie d'ordre supérieur et les groupoïdes multiples, le groupoïde de Teichmüller, son enseignement à Bangor. Correspondance d'un tiers, non transcrite (RIGHTS.md).}

\page{283}
\note{Suite de la lettre de R.~Brown du 16.4.1991 ; non transcrite (RIGHTS.md).}

\page{284}
\note{Suite de la lettre de R.~Brown du 16.4.1991 ; non transcrite (RIGHTS.md).}

\page{285}
\note{Suite de la lettre de R.~Brown du 16.4.1991 ; non transcrite (RIGHTS.md).}

\page{286}
\note{Suite de la lettre de R.~Brown du 16.4.1991 ; non transcrite (RIGHTS.md).}

\section*{Résumé de cours de J.-P. Serre, 1990--1991 (pages 288 à 300)}

\page{288}
\note{Début d'un document dactylographié de J.-P. Serre : son résumé de cours et travaux au Collège de France pour 1990--1991 (pp.~288--301 de la cote ; la fin est au lot 16). Texte d'un tiers, déjà imprimé dans l'\emph{Annuaire du Collège de France} et dans Serre, \emph{Résumés des cours au Collège de France 1956--2000} (SMF, 2001) : non transcrit.}

\page{289}
\note{Suite du résumé de J.-P. Serre (1990--1991) ; non transcrit.}

\page{290}
\note{Suite du résumé de J.-P. Serre (1990--1991) ; non transcrit.}

\page{291}
\note{Suite du résumé de J.-P. Serre (1990--1991) ; non transcrit.}

\page{292}
\note{Suite du résumé de J.-P. Serre (1990--1991) ; non transcrit.}

\page{293}
\note{Suite du résumé de J.-P. Serre (1990--1991) ; non transcrit.}

\page{294}
\note{Suite du résumé de J.-P. Serre (1990--1991) ; non transcrit.}

\page{295}
\note{Suite du résumé de J.-P. Serre (1990--1991) ; non transcrit.}

\page{296}
\note{Suite du résumé de J.-P. Serre (1990--1991) ; non transcrit.}

\page{297}
\note{Suite du résumé de J.-P. Serre (1990--1991) ; non transcrit.}

\page{298}
\note{Suite du résumé de J.-P. Serre (1990--1991) ; non transcrit.}

\page{299}
\note{Suite du résumé de J.-P. Serre (1990--1991) ; non transcrit.}

\page{300}
\note{Suite du résumé de J.-P. Serre (1990--1991) ; non transcrit.}

\end{document}
